\documentclass[pdftex,francais]{beamer}

% Copyright 2004 by Till Tantau <tantau@users.sourceforge.net>.
%
% This file can be redistributed and/or modified under
% the terms of the GNU Public License, version 2.

%% \ifx\themename\undefined
%%   \def\themename{default}
%% \fi

\usetheme{lama}
%\usetheme{Madrid}
%\usecolortheme{crane}

\usepackage{multirow}
\usepackage[latin1]{inputenc}           %%%  
\usepackage[T1]{fontenc}                %%%
\usepackage[francais]{babel}            %%%

\usepackage{multimedia}
\usepackage{hyperref}
\usepackage{tikz}
%\newtheorem{theorem}{Th�or�me}

%\setbeamercovered{transparent}

\title[DGtal library]{DGtal: Digital Geometry Tools and Algorithms\\
     \url{http://liris.cnrs.fr/dgtal}
}

\author[The DGtal Team]{The DGtal Team}

\date{DGCI demo session, April 2011}


\graphicspath{{Figures/},{Images/},{Graphs/}}

%%% \AtBeginSection[]
%%% {
%%%   \begin{frame}<beamer>
%%%     \frametitle{Plan}
%%%     \tableofcontents[currentsection] %,currentsubsection]
%%%   \end{frame}
%%% }


\renewcommand{\vec}[1]{\mathbf{#1}}

% space of the real numbers
\newcommand{\R}{\ensuremath{\mathbb{R}}}
% space of the integer numbers
\newcommand{\Z}{\ensuremath{\mathbb{Z}}}
% Digitization process of step h (1)
\newcommand{\Dig}[1]{\ensuremath{\mathrm{Dig}_{#1}}}
% Family of shape
\newcommand{\SF}[0]{\ensuremath{\mathbb{F}}}
% Topological boundary of X (1).
\newcommand{\TB}[1]{\ensuremath{\partial #1}}
% Discrete geometric estimator of G (1)
\newcommand{\DGE}[1]{\ensuremath{E_{#1}}}
% Reference shape of digital object O (1) with grid step $h$ (2).
\newcommand{\RS}[2]{\ensuremath{R_{#1,#2}}}

% Digital contour.
\newcommand{\DC}{\ensuremath{C}}
% Continuous contour.
\newcommand{\CC}{\ensuremath{\mathcal{C}}}
% A point indexed by i (1) on the digital contour.
\newcommand{\PT}[1]{\ensuremath{\DC_{#1}}}
% A sequence of points indexed by i (1) on the digital contour.
\newcommand{\PTS}[2]{\ensuremath{\DC_{#1,#2}}}
% Predicate stating that the digital curve is a segment between indices 1 and 2
\newcommand{\SPRED}[2]{\ensuremath{S(#1,#2)}}
% ET logique
\newcommand{\AND}[0]{\ensuremath{\wedge}}
% OU logique
\newcommand{\OR}[0]{\ensuremath{\vee}}

% Tangent direction mapping of curve C (1)
\newcommand{\TGT}[1]{\ensuremath{\theta_{#1}}}
% Integral of squared curvature of curve C (1)
\newcommand{\ISC}[1]{\ensuremath{J[#1]}}
% Smallest possible tangent direction at constraint l (1)
\newcommand{\MinTD}[1]{\ensuremath{a_{#1}}}
% Largest possible tangent direction at constraint l (1)
\newcommand{\MaxTD}[1]{\ensuremath{b_{#1}}}
% Unknown tangent direction at constraint l (1)
\newcommand{\UnkTD}[1]{\ensuremath{t_{#1}}}

% ideal multiscale criterion
\newcommand{\IMSC}[2]{\ensuremath{\mu_{#1}(#2)}}
% multiscale profile
\newcommand{\MSP}[2]{\ensuremath{\mathcal{P}_{#1}(#2)}}
% threshold flat/curve
\newcommand{\CThreshold}[0]{\ensuremath{t_{f/c}}}
% threshold noise
\newcommand{\NThreshold}[0]{\ensuremath{t_{m}}}
% noise level
\newcommand{\NL}[0]{\ensuremath{\nu}}
% slope linear regression 
\newcommand{\SLR}[2]{\ensuremath{\theta_{#1}}}





% Pencil of maximal segments around a point.
\newcommand{\PE}[1]{\ensuremath{{\mathcal P}(#1)}}
% Tangent orientation of the maximal segment.
\newcommand{\TO}[1]{\ensuremath{\theta_#1}}
% lambda function for interpolation.
\newcommand{\LF}{\ensuremath{\lambda}}
% \lambda-MS tangent orientation.
\newcommand{\LTO}[1]{\ensuremath{\hat{\theta}(#1)}}
% \lambda-MS tangent orientation variation.
\newcommand{\LTOP}[1]{\ensuremath{\hat{\theta}'(#1)}}

%\definecolor{rougeSW_}{rgb}{0.968627,0.011765,0.015686}
\definecolor{darkgreen}{rgb}{0.0,0.6,0.0}
\definecolor{lightblue}{rgb}{0.5,0.5,1.0}
\definecolor{magenta}{rgb}{1.0,0.0,1.0}
\newcommand{\alertred}[1]{{\color{red}#1}}
\newcommand{\Cb}[1]{{\color{blue}#1}}
\newcommand{\Cdg}[1]{{\color{darkgreen}#1}}
\newcommand{\textmagenta}[1]{{\color{magenta}#1}}
\newcommand{\Implies}{{\ensuremath{\Rightarrow}}}

\newcommand{\Refs}[1]{{\color{lightblue}#1}}
\newcommand{\Cite}[1]{\Refs{[#1]}}
\newcommand{\Etal}{{\em et al.}}
\newtheorem{remark}{Remarque}
% Digitization process of step h (1)
\newcommand{\DigGh}[2]{\ensuremath{\mathrm{Dig}_{#2}(#1)}}
\newcommand{\BigT}{\ensuremath{\Theta}}
\newcommand{\BigO}{\ensuremath{O}}

\newcommand{\TAN}[0]{\ensuremath{\theta}}
% Position estimator
\newcommand{\EPOS}[0]{\ensuremath{\hat{{x}}}}
% Convexity Position estimator 
\newcommand{\ECONVPOS}[0]{\ensuremath{\hat{x}^\mathrm{conv}}}
% tangent estimator base MS.
\newcommand{\ETANMS}[0]{\ensuremath{\hat{\TAN}^{\text{MS}}}}
% tangent estimator base arete du CDP.
\newcommand{\ETANEDGE}[0]{\ensuremath{\hat{\TAN}^{\text{conv}}}}
% estimateur de longueur elementaire d'un surfel (1) sur le bord discretise de X(2) de pas h(3).


\setbeamercolor{qcolorb}{fg={blue!20!black},bg={blue!15!white}}
\setbeamercolor{qcolorub}{fg={blue!10!black},bg={blue!30!white}}
\setbeamercolor{qcolorlb}{fg={blue!20!black},bg={blue!8!white}}
\setbeamercolor{qcolorulb}{fg={blue!10!black},bg={blue!40!white}}
\setbeamercolor{qcolorg}{fg={green!20!black},bg={green!15!white}}
\setbeamercolor{qcolorug}{fg={green!10!black},bg={green!30!white}}
\setbeamercolor{qcolorlg}{fg={green!20!black},bg={green!8!white}}
\setbeamercolor{qcolorulg}{fg={green!10!black},bg={green!40!white}}
\setbeamercolor{qcolorr}{fg={red!20!black},bg={red!15!white}}
\setbeamercolor{qcolorur}{fg={red!10!black},bg={red!30!white}}
\setbeamercolor{qcolorlr}{fg={red!20!black},bg={red!8!white}}
\setbeamercolor{qcolorulr}{fg={red!10!black},bg={red!40!white}}
\newenvironment{myblockbluish}[2]%
	       {\begin{beamerboxesrounded}[lower=qcolorb,upper=qcolorub,width=#1,shadow=true]{#2}}{\end{beamerboxesrounded}}
\newenvironment{myblocklbluish}[2]%
	       {\begin{beamerboxesrounded}[lower=qcolorlb,upper=qcolorulb,width=#1,shadow=true]{#2}}{\end{beamerboxesrounded}}
\newenvironment{myblockgreenish}[2]%
	       {\begin{beamerboxesrounded}[lower=qcolorg,upper=qcolorug,width=#1,shadow=true]{#2}}{\end{beamerboxesrounded}}
\newenvironment{myblocklgreenish}[2]%
	       {\begin{beamerboxesrounded}[lower=qcolorlg,upper=qcolorulg,width=#1,shadow=true]{#2}}{\end{beamerboxesrounded}}
\newenvironment{myblockredish}[2]%
	       {\begin{beamerboxesrounded}[lower=qcolorr,upper=qcolorur,width=#1
,shadow=true]{#2}}{\end{beamerboxesrounded}}
\newenvironment{myblocklredish}[2]%
	       {\begin{beamerboxesrounded}[lower=qcolorlr,upper=qcolorulr,width=#1,shadow=true]{#2}}{\end{beamerboxesrounded}}

%%%%%%%%%%%%%%%%%%%%%%%%%%%%%%%%%%%%%%%%%%%%%%%%%%%%%%%%%%%%%%%%%%%%%%%%%%%%%%%
%%%%%%%%%%%%%%%%%%%%%%%%%%%%%%%%%%%%%%%%%%%%%%%%%%%%%%%%%%%%%%%%%%%%%%%%%%%%%%%
%%%%%%%%%%%%%%%%%%%%%%%%%%%%%%%%%%%%%%%%%%%%%%%%%%%%%%%%%%%%%%%%%%%%%%%%%%%%%%%
\begin{document}

\newlength{\unquart}
\setlength{\unquart}{0.21\textwidth}

%------------------------------------------------------------------------------
\begin{frame}
  \titlepage
\end{frame}
%------------------------------------------------------------------------------

%------------------------------------------------------------------------------
\begin{frame}%[allowframebreaks]
  \frametitle{DGtal: why, who}
  
  \small
  \begin{block}{Objectives}
    \small
    \begin{itemize}
    \item to make easier discrete geometry for the neophyte (student, researcher from another field, \ldots)
    \item to test quickly new ideas, with objective comparison wrt
      existant works
    \item to make easier the implementation of demonstrators
    \item to help spread our research results to other domains
    \item to pursue a federative project
    \end{itemize}
  \end{block}

  \small
  \begin{block}{Who ? for now \ldots}
    \small \vspace{-0.4cm}
    \begin{columns}
      \begin{column}{0.45\textwidth}
        \begin{itemize}
        \item LIRIS (Lyon)      
        \item Gipsa-lab (Grenoble)
        \item GREYC (Caen)
        \end{itemize}
      \end{column}
      \begin{column}{0.45\textwidth}
        \begin{itemize}
        \item LAMA (Chamb�ry)
        \item LORIA (Nancy)
        \end{itemize}
      \end{column}
    \end{columns}
  \end{block}
\end{frame}
%------------------------------------------------------------------------------

%------------------------------------------------------------------------------
\begin{frame}%[allowframebreaks]
  \frametitle{DGtal: what for ?}
  
  \small
  \begin{block}{Main features}
    \small
    \begin{itemize}
    \item to define digital objects in arbitrary dimension
    \item to propose algorithms for topological and geometric analysis
    \item to provide I/O mechanisms and visualization tools
    \end{itemize}
  \end{block}
  \medskip

  \begin{tabular}{cccc}
    \includegraphics[height=0.15\textheight]{exampleDSS-3}
    &
    \includegraphics[height=0.2\textheight]{edt-2d}
    &
    \includegraphics[height=0.2\textheight]{object-3d-18-6}
    &
    \includegraphics[height=0.2\textwidth]{thinning-3d}
    \\
    DSS & DT & Objects & Thinning\\
  \end{tabular}
\end{frame}
%------------------------------------------------------------------------------
%------------------------------------------------------------------------------
\begin{frame}%[allowframebreaks]
  \frametitle{DGtal philosophy and structure}

  \begin{block}{}
    \small
    - Genericity and efficiency: C++ library, concepts\\
    - LGPL or GPL with restrictions\\
    - user friendly, not necessarily kernel-developer friendly
  \end{block}
  \begin{myblockbluish}{\textwidth}{\centering Kernel}
    \begin{columns}
      \begin{column}{0.44\textwidth}
        \begin{myblocklbluish}{\textwidth}{Basic types, data structures}
          \scriptsize
          - digital space, point, sets, lin. algebra\\
          - software infrastructure: trace, concept validation, \ldots
        \end{myblocklbluish}
      \end{column}
      \begin{column}{0.44\textwidth}
        \begin{myblocklbluish}{\textwidth}{Images}
          \scriptsize
          - generic container\\
          - several implementation: standard, other adapted to big images
        \end{myblocklbluish}
      \end{column}
    \end{columns}
  \end{myblockbluish}
  \begin{myblockgreenish}{\textwidth}{\centering Base modules}
   \begin{columns}
      \begin{column}{0.44\textwidth}
        \begin{myblocklgreenish}{\textwidth}{Topology module}
          \scriptsize
          - digital topology: connectedness, border, simple points\\
          - grid topology: cells, contours, surfaces, tracking
        \end{myblocklgreenish}
      \end{column}
      \begin{column}{0.44\textwidth}
        \begin{myblocklgreenish}{\textwidth}{Geometry module}
          \scriptsize
          - primitives, DSS recognition\\
          - contour analysis: decomposition, convexity\\
          - volumetric analysis: distance transforms, medial axis
        \end{myblocklgreenish}
      \end{column}
    \end{columns}
   \end{myblockgreenish}

\end{frame}
%------------------------------------------------------------------------------

%------------------------------------------------------------------------------
\begin{frame}%[allowframebreaks]
  \frametitle{Other modules, or modules with external dependencies}
  
  \begin{columns}
    \scriptsize
    \begin{column}{0.32\textwidth}
      \begin{myblocklredish}{\textwidth}{2D vector export (SVG,\ldots)}
        \includegraphics[width=\textwidth]{house-layers4-4}
      \end{myblocklredish}
    \end{column}
    \begin{column}{0.32\textwidth}
      \begin{myblocklredish}{\textwidth}{3D visualization}
        \includegraphics[width=\textwidth]{al-clipped}
      \end{myblocklredish}
    \end{column}
    \begin{column}{0.32\textwidth}
      \begin{myblocklredish}{\textwidth}{Backends}
        \begin{itemize}
        \item Kiteware's ITK
        \item VIGRA (soon)
        \end{itemize}
      \end{myblocklredish}\\
      \begin{myblocklredish}{\textwidth}{Import/export}
        \begin{itemize}
        \item images (ImageMagick)
        \item volumes (libvol)
        \end{itemize}
      \end{myblocklredish}
      \begin{myblocklredish}{\textwidth}{Project modules}
        \begin{itemize}
        \item noisy objects (GeoDIB)
        \item \ldots
        \end{itemize}
      \end{myblocklredish}
    \end{column}
  \end{columns}

\end{frame}
%------------------------------------------------------------------------------

%------------------------------------------------------------------------------
\begin{frame}%[allowframebreaks]
  \frametitle{DGtal Roadmap}
\small
  \begin{block}{New in milestone 0.3 (now)}
    \small
    \begin{itemize}
    \item Kernel update: integer genericity, concept checking, domain iterators
    \item 2D, nD-volumetric geometry: primitive decomposition, tangential cover, reverse distance transformation
    \item 3D visualization with QGLviewer: stream mechanism as in 2D
    \item grid or interpixel topology: cells, digital surfaces,
      surface tracking
    \item shape construction
    \end{itemize}
  \end{block}

  \begin{myblockredish}{\textwidth}{Join DGtal}
    \small
    \begin{itemize}
    \item new contributors are welcome (new bug-reporters, documentation readers are welcome too)
    \item collaborative forge, development infrastructure
    \item DGtal week this summer
    \end{itemize}
  \end{myblockredish}
  
\end{frame}

%------------------------------------------------------------------------------
%------------------------------------------------------------------------------
\begin{frame}
  \frametitle{DGtal Team}

   \begin{center}
     \includegraphics[width=0.3\textwidth]{dgtal-logo}\\
     \url{http://liris.cnrs.fr/dgtal}
   \end{center}

   %%       \begin{tikzpicture}[remember picture,overlay]
   %%         \begin{scope}[scale=0.1]
  %\definecolor{lamaColor}{rgb}{0.227451,0.48039,0.035294}
  \pgfsetfillcolor{univColor}
  \pgfsetstrokecolor{univColor}
  \pgfpathmoveto{\pgfqpoint{12.871cm}{17.371cm}}
\pgfpathcurveto{\pgfqpoint{13.047cm}{17.502cm}}{\pgfqpoint{13.284cm}{17.544cm}}{\pgfqpoint{13.492cm}{17.472cm}}
\pgfpathcurveto{\pgfqpoint{13.355cm}{17.338cm}}{\pgfqpoint{13.222cm}{17.192cm}}{\pgfqpoint{13.154cm}{17.009cm}}
\pgfpathcurveto{\pgfqpoint{13.073cm}{16.82cm}}{\pgfqpoint{13.072cm}{16.607cm}}{\pgfqpoint{13.12cm}{16.41cm}}
\pgfpathcurveto{\pgfqpoint{13.279cm}{16.41cm}}{\pgfqpoint{13.438cm}{16.41cm}}{\pgfqpoint{13.597cm}{16.412cm}}
\pgfpathcurveto{\pgfqpoint{13.933cm}{16.84cm}}{\pgfqpoint{14.362cm}{17.257cm}}{\pgfqpoint{14.924cm}{17.34cm}}
\pgfpathcurveto{\pgfqpoint{14.845cm}{17.059cm}}{\pgfqpoint{14.576cm}{16.901cm}}{\pgfqpoint{14.453cm}{16.645cm}}
\pgfpathcurveto{\pgfqpoint{14.364cm}{16.477cm}}{\pgfqpoint{14.274cm}{16.311cm}}{\pgfqpoint{14.183cm}{16.144cm}}
\pgfpathcurveto{\pgfqpoint{14.289cm}{16.046cm}}{\pgfqpoint{14.402cm}{15.956cm}}{\pgfqpoint{14.512cm}{15.863cm}}
\pgfpathcurveto{\pgfqpoint{14.674cm}{15.729cm}}{\pgfqpoint{14.785cm}{15.548cm}}{\pgfqpoint{14.938cm}{15.406cm}}
\pgfpathcurveto{\pgfqpoint{15.057cm}{15.281cm}}{\pgfqpoint{15.198cm}{15.172cm}}{\pgfqpoint{15.288cm}{15.023cm}}
\pgfpathcurveto{\pgfqpoint{15.367cm}{14.824cm}}{\pgfqpoint{15.419cm}{14.613cm}}{\pgfqpoint{15.451cm}{14.402cm}}
\pgfpathcurveto{\pgfqpoint{15.457cm}{14.139cm}}{\pgfqpoint{15.213cm}{13.978cm}}{\pgfqpoint{15.06cm}{13.798cm}}
\pgfpathcurveto{\pgfqpoint{14.863cm}{13.783cm}}{\pgfqpoint{14.676cm}{13.85cm}}{\pgfqpoint{14.49cm}{13.902cm}}
\pgfpathcurveto{\pgfqpoint{14.23cm}{13.982cm}}{\pgfqpoint{13.97cm}{13.855cm}}{\pgfqpoint{13.708cm}{13.856cm}}
\pgfpathcurveto{\pgfqpoint{13.616cm}{13.732cm}}{\pgfqpoint{13.533cm}{13.601cm}}{\pgfqpoint{13.478cm}{13.457cm}}
\pgfpathcurveto{\pgfqpoint{13.382cm}{13.225cm}}{\pgfqpoint{13.353cm}{12.973cm}}{\pgfqpoint{13.272cm}{12.736cm}}
\pgfpathcurveto{\pgfqpoint{13.211cm}{12.529cm}}{\pgfqpoint{13.12cm}{12.329cm}}{\pgfqpoint{13.097cm}{12.114cm}}
\pgfpathcurveto{\pgfqpoint{13.021cm}{11.535cm}}{\pgfqpoint{12.974cm}{10.954cm}}{\pgfqpoint{12.895cm}{10.376cm}}
\pgfpathcurveto{\pgfqpoint{12.866cm}{10.205cm}}{\pgfqpoint{12.85cm}{10.032cm}}{\pgfqpoint{12.833cm}{9.859cm}}
\pgfpathcurveto{\pgfqpoint{12.771cm}{9.522cm}}{\pgfqpoint{12.549cm}{9.247cm}}{\pgfqpoint{12.418cm}{8.937cm}}
\pgfpathcurveto{\pgfqpoint{12.263cm}{8.645cm}}{\pgfqpoint{12.298cm}{8.298cm}}{\pgfqpoint{12.159cm}{8cm}}
\pgfpathcurveto{\pgfqpoint{12.093cm}{7.862cm}}{\pgfqpoint{12.071cm}{7.708cm}}{\pgfqpoint{12.056cm}{7.557cm}}
\pgfpathcurveto{\pgfqpoint{12.037cm}{7.377cm}}{\pgfqpoint{11.969cm}{7.207cm}}{\pgfqpoint{11.934cm}{7.031cm}}
\pgfpathcurveto{\pgfqpoint{11.895cm}{6.868cm}}{\pgfqpoint{11.876cm}{6.7cm}}{\pgfqpoint{11.816cm}{6.544cm}}
\pgfpathcurveto{\pgfqpoint{11.738cm}{6.367cm}}{\pgfqpoint{11.558cm}{6.275cm}}{\pgfqpoint{11.412cm}{6.163cm}}
\pgfpathcurveto{\pgfqpoint{11.281cm}{6.071cm}}{\pgfqpoint{11.195cm}{5.933cm}}{\pgfqpoint{11.084cm}{5.82cm}}
\pgfpathcurveto{\pgfqpoint{11.015cm}{5.749cm}}{\pgfqpoint{10.949cm}{5.673cm}}{\pgfqpoint{10.915cm}{5.578cm}}
\pgfpathcurveto{\pgfqpoint{10.774cm}{5.205cm}}{\pgfqpoint{10.628cm}{4.832cm}}{\pgfqpoint{10.452cm}{4.474cm}}
\pgfpathcurveto{\pgfqpoint{10.296cm}{4.16cm}}{\pgfqpoint{10.014cm}{3.927cm}}{\pgfqpoint{9.874cm}{3.605cm}}
\pgfpathcurveto{\pgfqpoint{9.777cm}{3.384cm}}{\pgfqpoint{9.838cm}{3.133cm}}{\pgfqpoint{9.757cm}{2.907cm}}
\pgfpathcurveto{\pgfqpoint{9.673cm}{2.659cm}}{\pgfqpoint{9.55cm}{2.422cm}}{\pgfqpoint{9.501cm}{2.162cm}}
\pgfpathcurveto{\pgfqpoint{9.438cm}{1.887cm}}{\pgfqpoint{9.466cm}{1.6cm}}{\pgfqpoint{9.411cm}{1.323cm}}
\pgfpathcurveto{\pgfqpoint{9.359cm}{0.911cm}}{\pgfqpoint{9.159cm}{0.127cm}}{\pgfqpoint{9.157cm}{-0.287cm}}
\pgfpathlineto{\pgfqpoint{8.598cm}{-0.287cm}}
\pgfpathcurveto{\pgfqpoint{8.636cm}{-0.112cm}}{\pgfqpoint{8.616cm}{0.431cm}}{\pgfqpoint{8.606cm}{0.611cm}}
\pgfpathcurveto{\pgfqpoint{8.574cm}{1.042cm}}{\pgfqpoint{8.741cm}{1.453cm}}{\pgfqpoint{8.753cm}{1.881cm}}
\pgfpathcurveto{\pgfqpoint{8.76cm}{2.133cm}}{\pgfqpoint{8.895cm}{2.431cm}}{\pgfqpoint{8.689cm}{2.64cm}}
\pgfpathcurveto{\pgfqpoint{8.734cm}{2.756cm}}{\pgfqpoint{8.718cm}{2.878cm}}{\pgfqpoint{8.678cm}{2.992cm}}
\pgfpathcurveto{\pgfqpoint{8.801cm}{3.161cm}}{\pgfqpoint{8.782cm}{3.377cm}}{\pgfqpoint{8.697cm}{3.557cm}}
\pgfpathcurveto{\pgfqpoint{8.734cm}{3.569cm}}{\pgfqpoint{8.771cm}{3.581cm}}{\pgfqpoint{8.809cm}{3.593cm}}
\pgfpathcurveto{\pgfqpoint{8.785cm}{3.751cm}}{\pgfqpoint{8.817cm}{3.907cm}}{\pgfqpoint{8.838cm}{4.062cm}}
\pgfpathcurveto{\pgfqpoint{8.797cm}{4.27cm}}{\pgfqpoint{8.75cm}{4.487cm}}{\pgfqpoint{8.821cm}{4.694cm}}
\pgfpathcurveto{\pgfqpoint{8.777cm}{4.696cm}}{\pgfqpoint{8.733cm}{4.697cm}}{\pgfqpoint{8.689cm}{4.699cm}}
\pgfpathcurveto{\pgfqpoint{8.616cm}{4.149cm}}{\pgfqpoint{8.609cm}{3.588cm}}{\pgfqpoint{8.48cm}{3.046cm}}
\pgfpathcurveto{\pgfqpoint{8.448cm}{2.864cm}}{\pgfqpoint{8.316cm}{2.726cm}}{\pgfqpoint{8.247cm}{2.559cm}}
\pgfpathcurveto{\pgfqpoint{7.999cm}{2.012cm}}{\pgfqpoint{7.94cm}{1.405cm}}{\pgfqpoint{7.859cm}{0.816cm}}
\pgfpathcurveto{\pgfqpoint{7.904cm}{0.566cm}}{\pgfqpoint{7.716cm}{-0.037cm}}{\pgfqpoint{7.766cm}{-0.287cm}}
\pgfpathlineto{\pgfqpoint{7.062cm}{-0.293cm}}
\pgfpathcurveto{\pgfqpoint{7.117cm}{-0.056cm}}{\pgfqpoint{7.137cm}{0.497cm}}{\pgfqpoint{7.187cm}{0.758cm}}
\pgfpathcurveto{\pgfqpoint{7.205cm}{1.019cm}}{\pgfqpoint{7.112cm}{1.316cm}}{\pgfqpoint{7.282cm}{1.546cm}}
\pgfpathcurveto{\pgfqpoint{7.416cm}{1.71cm}}{\pgfqpoint{7.421cm}{1.927cm}}{\pgfqpoint{7.467cm}{2.124cm}}
\pgfpathcurveto{\pgfqpoint{7.502cm}{2.362cm}}{\pgfqpoint{7.364cm}{2.592cm}}{\pgfqpoint{7.417cm}{2.833cm}}
\pgfpathcurveto{\pgfqpoint{7.438cm}{3.186cm}}{\pgfqpoint{7.407cm}{3.539cm}}{\pgfqpoint{7.371cm}{3.89cm}}
\pgfpathcurveto{\pgfqpoint{7.353cm}{4.026cm}}{\pgfqpoint{7.37cm}{4.162cm}}{\pgfqpoint{7.397cm}{4.295cm}}
\pgfpathcurveto{\pgfqpoint{7.346cm}{4.338cm}}{\pgfqpoint{7.282cm}{4.373cm}}{\pgfqpoint{7.253cm}{4.436cm}}
\pgfpathcurveto{\pgfqpoint{7.218cm}{4.513cm}}{\pgfqpoint{7.244cm}{4.599cm}}{\pgfqpoint{7.251cm}{4.678cm}}
\pgfpathcurveto{\pgfqpoint{7.121cm}{4.753cm}}{\pgfqpoint{6.979cm}{4.806cm}}{\pgfqpoint{6.854cm}{4.888cm}}
\pgfpathcurveto{\pgfqpoint{6.691cm}{5.032cm}}{\pgfqpoint{6.564cm}{5.215cm}}{\pgfqpoint{6.384cm}{5.341cm}}
\pgfpathcurveto{\pgfqpoint{6.299cm}{5.407cm}}{\pgfqpoint{6.184cm}{5.444cm}}{\pgfqpoint{6.132cm}{5.546cm}}
\pgfpathcurveto{\pgfqpoint{6.024cm}{5.734cm}}{\pgfqpoint{5.921cm}{5.936cm}}{\pgfqpoint{5.73cm}{6.054cm}}
\pgfpathcurveto{\pgfqpoint{5.514cm}{5.837cm}}{\pgfqpoint{5.538cm}{5.513cm}}{\pgfqpoint{5.511cm}{5.232cm}}
\pgfpathcurveto{\pgfqpoint{5.517cm}{4.945cm}}{\pgfqpoint{5.435cm}{4.666cm}}{\pgfqpoint{5.423cm}{4.381cm}}
\pgfpathcurveto{\pgfqpoint{5.347cm}{4.334cm}}{\pgfqpoint{5.258cm}{4.292cm}}{\pgfqpoint{5.22cm}{4.205cm}}
\pgfpathcurveto{\pgfqpoint{5.075cm}{3.928cm}}{\pgfqpoint{4.957cm}{3.636cm}}{\pgfqpoint{4.784cm}{3.374cm}}
\pgfpathcurveto{\pgfqpoint{4.69cm}{3.24cm}}{\pgfqpoint{4.762cm}{3.058cm}}{\pgfqpoint{4.671cm}{2.922cm}}
\pgfpathcurveto{\pgfqpoint{4.576cm}{2.772cm}}{\pgfqpoint{4.528cm}{2.598cm}}{\pgfqpoint{4.508cm}{2.423cm}}
\pgfpathcurveto{\pgfqpoint{4.494cm}{2.284cm}}{\pgfqpoint{4.447cm}{2.142cm}}{\pgfqpoint{4.485cm}{2.003cm}}
\pgfpathcurveto{\pgfqpoint{4.545cm}{1.724cm}}{\pgfqpoint{4.484cm}{1.434cm}}{\pgfqpoint{4.562cm}{1.157cm}}
\pgfpathcurveto{\pgfqpoint{4.608cm}{0.881cm}}{\pgfqpoint{4.747cm}{0.633cm}}{\pgfqpoint{4.806cm}{0.361cm}}
\pgfpathcurveto{\pgfqpoint{4.826cm}{0.243cm}}{\pgfqpoint{4.969cm}{-0.203cm}}{\pgfqpoint{5.044cm}{-0.287cm}}
\pgfpathlineto{\pgfqpoint{4.335cm}{-0.287cm}}
\pgfpathcurveto{\pgfqpoint{4.321cm}{-0.071cm}}{\pgfqpoint{4.039cm}{0.473cm}}{\pgfqpoint{4.01cm}{0.685cm}}
\pgfpathcurveto{\pgfqpoint{3.985cm}{0.818cm}}{\pgfqpoint{3.881cm}{0.926cm}}{\pgfqpoint{3.881cm}{1.067cm}}
\pgfpathcurveto{\pgfqpoint{3.869cm}{1.272cm}}{\pgfqpoint{3.793cm}{1.464cm}}{\pgfqpoint{3.729cm}{1.656cm}}
\pgfpathcurveto{\pgfqpoint{3.679cm}{1.635cm}}{\pgfqpoint{3.628cm}{1.614cm}}{\pgfqpoint{3.578cm}{1.593cm}}
\pgfpathcurveto{\pgfqpoint{3.603cm}{1.7cm}}{\pgfqpoint{3.633cm}{1.805cm}}{\pgfqpoint{3.667cm}{1.91cm}}
\pgfpathcurveto{\pgfqpoint{3.629cm}{1.93cm}}{\pgfqpoint{3.591cm}{1.95cm}}{\pgfqpoint{3.553cm}{1.97cm}}
\pgfpathcurveto{\pgfqpoint{3.516cm}{2.225cm}}{\pgfqpoint{3.378cm}{2.445cm}}{\pgfqpoint{3.269cm}{2.673cm}}
\pgfpathcurveto{\pgfqpoint{3.182cm}{2.85cm}}{\pgfqpoint{3.343cm}{3.013cm}}{\pgfqpoint{3.374cm}{3.184cm}}
\pgfpathcurveto{\pgfqpoint{3.42cm}{3.421cm}}{\pgfqpoint{3.592cm}{3.615cm}}{\pgfqpoint{3.607cm}{3.86cm}}
\pgfpathcurveto{\pgfqpoint{3.634cm}{4.129cm}}{\pgfqpoint{3.61cm}{4.4cm}}{\pgfqpoint{3.583cm}{4.667cm}}
\pgfpathcurveto{\pgfqpoint{3.538cm}{4.739cm}}{\pgfqpoint{3.467cm}{4.789cm}}{\pgfqpoint{3.408cm}{4.848cm}}
\pgfpathcurveto{\pgfqpoint{3.34cm}{4.735cm}}{\pgfqpoint{3.254cm}{4.631cm}}{\pgfqpoint{3.201cm}{4.509cm}}
\pgfpathcurveto{\pgfqpoint{3.161cm}{4.337cm}}{\pgfqpoint{3.133cm}{4.162cm}}{\pgfqpoint{3.087cm}{3.992cm}}
\pgfpathcurveto{\pgfqpoint{3.034cm}{3.749cm}}{\pgfqpoint{2.957cm}{3.505cm}}{\pgfqpoint{2.981cm}{3.254cm}}
\pgfpathcurveto{\pgfqpoint{2.998cm}{3.05cm}}{\pgfqpoint{2.911cm}{2.858cm}}{\pgfqpoint{2.897cm}{2.657cm}}
\pgfpathcurveto{\pgfqpoint{2.844cm}{2.089cm}}{\pgfqpoint{2.623cm}{1.298cm}}{\pgfqpoint{2.542cm}{0.736cm}}
\pgfpathcurveto{\pgfqpoint{2.535cm}{0.584cm}}{\pgfqpoint{2.487cm}{0.074cm}}{\pgfqpoint{2.526cm}{-0.073cm}}
\pgfpathcurveto{\pgfqpoint{2.566cm}{-0.224cm}}{\pgfqpoint{2.548cm}{-0.189cm}}{\pgfqpoint{2.601cm}{-0.3cm}}
\pgfpathlineto{\pgfqpoint{1.787cm}{-0.3cm}}
\pgfpathcurveto{\pgfqpoint{1.814cm}{-0.193cm}}{\pgfqpoint{1.874cm}{-0.105cm}}{\pgfqpoint{1.889cm}{0.044cm}}
\pgfpathcurveto{\pgfqpoint{1.801cm}{0.264cm}}{\pgfqpoint{1.789cm}{0.869cm}}{\pgfqpoint{1.808cm}{1.106cm}}
\pgfpathcurveto{\pgfqpoint{1.838cm}{1.477cm}}{\pgfqpoint{1.791cm}{1.85cm}}{\pgfqpoint{1.717cm}{2.213cm}}
\pgfpathcurveto{\pgfqpoint{1.674cm}{2.427cm}}{\pgfqpoint{1.566cm}{2.62cm}}{\pgfqpoint{1.517cm}{2.832cm}}
\pgfpathcurveto{\pgfqpoint{1.431cm}{3.171cm}}{\pgfqpoint{1.491cm}{3.522cm}}{\pgfqpoint{1.544cm}{3.861cm}}
\pgfpathcurveto{\pgfqpoint{1.57cm}{4.123cm}}{\pgfqpoint{1.679cm}{4.369cm}}{\pgfqpoint{1.698cm}{4.632cm}}
\pgfpathcurveto{\pgfqpoint{1.72cm}{4.836cm}}{\pgfqpoint{1.688cm}{5.048cm}}{\pgfqpoint{1.586cm}{5.227cm}}
\pgfpathcurveto{\pgfqpoint{1.445cm}{5.468cm}}{\pgfqpoint{1.46cm}{5.749cm}}{\pgfqpoint{1.411cm}{6.015cm}}
\pgfpathcurveto{\pgfqpoint{1.359cm}{6.448cm}}{\pgfqpoint{1.316cm}{6.906cm}}{\pgfqpoint{1.236cm}{7.348cm}}
\pgfpathcurveto{\pgfqpoint{1.166cm}{7.741cm}}{\pgfqpoint{1.13cm}{8.12cm}}{\pgfqpoint{1.117cm}{8.518cm}}
\pgfpathcurveto{\pgfqpoint{1.098cm}{8.956cm}}{\pgfqpoint{1.092cm}{9.34cm}}{\pgfqpoint{1.055cm}{9.738cm}}
\pgfpathcurveto{\pgfqpoint{0.995cm}{10.28cm}}{\pgfqpoint{1.061cm}{10.411cm}}{\pgfqpoint{0.875cm}{10.58cm}}
\pgfpathcurveto{\pgfqpoint{0.599cm}{10.667cm}}{\pgfqpoint{0.379cm}{10.403cm}}{\pgfqpoint{0.223cm}{10.104cm}}
\pgfpathcurveto{\pgfqpoint{0.058cm}{9.787cm}}{\pgfqpoint{0.037cm}{9.667cm}}{\pgfqpoint{-0.293cm}{9.625cm}}
\pgfpathcurveto{\pgfqpoint{-0.649cm}{9.752cm}}{\pgfqpoint{-0.36cm}{10.138cm}}{\pgfqpoint{-0.278cm}{10.473cm}}
\pgfpathcurveto{\pgfqpoint{-0.162cm}{10.755cm}}{\pgfqpoint{-0.157cm}{10.742cm}}{\pgfqpoint{0.029cm}{11.021cm}}
\pgfpathcurveto{\pgfqpoint{0.203cm}{11.231cm}}{\pgfqpoint{0.297cm}{11.342cm}}{\pgfqpoint{0.51cm}{11.471cm}}
\pgfpathcurveto{\pgfqpoint{0.946cm}{11.711cm}}{\pgfqpoint{0.959cm}{11.712cm}}{\pgfqpoint{1.513cm}{11.858cm}}
\pgfpathcurveto{\pgfqpoint{2.418cm}{12.056cm}}{\pgfqpoint{3.456cm}{11.981cm}}{\pgfqpoint{4.16cm}{11.95cm}}
\pgfpathcurveto{\pgfqpoint{4.633cm}{11.992cm}}{\pgfqpoint{4.837cm}{12.018cm}}{\pgfqpoint{5.54cm}{11.987cm}}
\pgfpathcurveto{\pgfqpoint{6.164cm}{12.012cm}}{\pgfqpoint{6.805cm}{12.05cm}}{\pgfqpoint{7.409cm}{11.9cm}}
\pgfpathcurveto{\pgfqpoint{8.076cm}{11.751cm}}{\pgfqpoint{8.702cm}{11.467cm}}{\pgfqpoint{9.367cm}{11.313cm}}
\pgfpathcurveto{\pgfqpoint{9.811cm}{11.154cm}}{\pgfqpoint{10.037cm}{11.15cm}}{\pgfqpoint{10.371cm}{11.299cm}}
\pgfpathcurveto{\pgfqpoint{10.685cm}{11.615cm}}{\pgfqpoint{10.702cm}{11.578cm}}{\pgfqpoint{10.817cm}{11.861cm}}
\pgfpathcurveto{\pgfqpoint{10.912cm}{12.201cm}}{\pgfqpoint{10.958cm}{12.379cm}}{\pgfqpoint{11.046cm}{12.714cm}}
\pgfpathcurveto{\pgfqpoint{11.082cm}{12.857cm}}{\pgfqpoint{11.137cm}{12.994cm}}{\pgfqpoint{11.158cm}{13.14cm}}
\pgfpathcurveto{\pgfqpoint{11.235cm}{13.649cm}}{\pgfqpoint{11.354cm}{14.15cm}}{\pgfqpoint{11.442cm}{14.657cm}}
\pgfpathcurveto{\pgfqpoint{11.544cm}{15.242cm}}{\pgfqpoint{11.768cm}{15.808cm}}{\pgfqpoint{12.117cm}{16.29cm}}
\pgfpathcurveto{\pgfqpoint{12.243cm}{16.471cm}}{\pgfqpoint{12.467cm}{16.556cm}}{\pgfqpoint{12.594cm}{16.736cm}}
\pgfpathcurveto{\pgfqpoint{12.673cm}{16.952cm}}{\pgfqpoint{12.728cm}{17.185cm}}{\pgfqpoint{12.871cm}{17.371cm}}
\pgfpathclose
\pgfusepath{fill}
%\draw[fill=white, draw=none](13.7cm,15.3cm)circle(0.4cm);
%\draw[fill=fg, draw=none](13.8cm,15.37cm)circle(0.15cm);



\definecolor{textLogoColor}{rgb}{0.4,0.4,0.4}
\pgfsetfillcolor{textLogoColor}
\pgfsetstrokecolor{textLogoColor}
%%L
\begin{scope}[xshift=-3cm]
\pgfpathmoveto{\pgfqpoint{2.144cm}{-1.377cm}}
\pgfpathlineto{\pgfqpoint{2.144cm}{-3.611cm}}
\pgfpathlineto{\pgfqpoint{3.655cm}{-3.611cm}}
\pgfpathcurveto{\pgfqpoint{3.776cm}{-3.611cm}}{\pgfqpoint{3.868cm}{-3.636cm}}{\pgfqpoint{3.931cm}{-3.685cm}}
\pgfpathcurveto{\pgfqpoint{3.996cm}{-3.733cm}}{\pgfqpoint{4.029cm}{-3.795cm}}{\pgfqpoint{4.029cm}{-3.869cm}}
\pgfpathcurveto{\pgfqpoint{4.029cm}{-3.944cm}}{\pgfqpoint{3.997cm}{-4.006cm}}{\pgfqpoint{3.934cm}{-4.053cm}}
\pgfpathcurveto{\pgfqpoint{3.87cm}{-4.099cm}}{\pgfqpoint{3.777cm}{-4.122cm}}{\pgfqpoint{3.655cm}{-4.122cm}}
\pgfpathlineto{\pgfqpoint{1.856cm}{-4.122cm}}
\pgfpathcurveto{\pgfqpoint{1.693cm}{-4.122cm}}{\pgfqpoint{1.576cm}{-4.092cm}}{\pgfqpoint{1.504cm}{-4.032cm}}
\pgfpathcurveto{\pgfqpoint{1.434cm}{-3.972cm}}{\pgfqpoint{1.399cm}{-3.875cm}}{\pgfqpoint{1.399cm}{-3.741cm}}
\pgfpathlineto{\pgfqpoint{1.399cm}{-1.377cm}}
\pgfpathcurveto{\pgfqpoint{1.399cm}{-1.252cm}}{\pgfqpoint{1.432cm}{-1.158cm}}{\pgfqpoint{1.499cm}{-1.095cm}}
\pgfpathcurveto{\pgfqpoint{1.568cm}{-1.032cm}}{\pgfqpoint{1.657cm}{-1.001cm}}{\pgfqpoint{1.768cm}{-1.001cm}}
\pgfpathcurveto{\pgfqpoint{1.88cm}{-1.001cm}}{\pgfqpoint{1.97cm}{-1.032cm}}{\pgfqpoint{2.039cm}{-1.095cm}}
\pgfpathcurveto{\pgfqpoint{2.109cm}{-1.156cm}}{\pgfqpoint{2.144cm}{-1.25cm}}{\pgfqpoint{2.144cm}{-1.377cm}}
\pgfusepath{fill}
\end{scope}

%% A
\begin{scope}[xshift=-1cm]
\pgfpathmoveto{\pgfqpoint{6.8cm}{-3.775cm}}
\pgfpathlineto{\pgfqpoint{6.624cm}{-3.389cm}}
\pgfpathlineto{\pgfqpoint{5.128cm}{-3.389cm}}
\pgfpathlineto{\pgfqpoint{4.953cm}{-3.783cm}}
\pgfpathcurveto{\pgfqpoint{4.884cm}{-3.937cm}}{\pgfqpoint{4.825cm}{-4.04cm}}{\pgfqpoint{4.777cm}{-4.095cm}}
\pgfpathcurveto{\pgfqpoint{4.728cm}{-4.148cm}}{\pgfqpoint{4.649cm}{-4.174cm}}{\pgfqpoint{4.538cm}{-4.174cm}}
\pgfpathcurveto{\pgfqpoint{4.445cm}{-4.174cm}}{\pgfqpoint{4.362cm}{-4.146cm}}{\pgfqpoint{4.29cm}{-4.089cm}}
\pgfpathcurveto{\pgfqpoint{4.218cm}{-4.031cm}}{\pgfqpoint{4.182cm}{-3.967cm}}{\pgfqpoint{4.182cm}{-3.894cm}}
\pgfpathcurveto{\pgfqpoint{4.182cm}{-3.852cm}}{\pgfqpoint{4.19cm}{-3.809cm}}{\pgfqpoint{4.207cm}{-3.764cm}}
\pgfpathcurveto{\pgfqpoint{4.224cm}{-3.719cm}}{\pgfqpoint{4.251cm}{-3.657cm}}{\pgfqpoint{4.29cm}{-3.578cm}}
\pgfpathlineto{\pgfqpoint{5.231cm}{-1.585cm}}
\pgfpathcurveto{\pgfqpoint{5.258cm}{-1.527cm}}{\pgfqpoint{5.29cm}{-1.459cm}}{\pgfqpoint{5.326cm}{-1.38cm}}
\pgfpathcurveto{\pgfqpoint{5.365cm}{-1.299cm}}{\pgfqpoint{5.405cm}{-1.232cm}}{\pgfqpoint{5.447cm}{-1.179cm}}
\pgfpathcurveto{\pgfqpoint{5.49cm}{-1.126cm}}{\pgfqpoint{5.546cm}{-1.083cm}}{\pgfqpoint{5.615cm}{-1.051cm}}
\pgfpathcurveto{\pgfqpoint{5.685cm}{-1.017cm}}{\pgfqpoint{5.771cm}{-1.001cm}}{\pgfqpoint{5.874cm}{-1.001cm}}
\pgfpathcurveto{\pgfqpoint{5.977cm}{-1.001cm}}{\pgfqpoint{6.063cm}{-1.017cm}}{\pgfqpoint{6.132cm}{-1.051cm}}
\pgfpathcurveto{\pgfqpoint{6.202cm}{-1.083cm}}{\pgfqpoint{6.258cm}{-1.125cm}}{\pgfqpoint{6.3cm}{-1.176cm}}
\pgfpathcurveto{\pgfqpoint{6.344cm}{-1.228cm}}{\pgfqpoint{6.38cm}{-1.284cm}}{\pgfqpoint{6.408cm}{-1.344cm}}
\pgfpathcurveto{\pgfqpoint{6.438cm}{-1.403cm}}{\pgfqpoint{6.476cm}{-1.481cm}}{\pgfqpoint{6.521cm}{-1.581cm}}
\pgfpathlineto{\pgfqpoint{7.482cm}{-3.561cm}}
\pgfpathcurveto{\pgfqpoint{7.557cm}{-3.712cm}}{\pgfqpoint{7.595cm}{-3.821cm}}{\pgfqpoint{7.595cm}{-3.89cm}}
\pgfpathcurveto{\pgfqpoint{7.595cm}{-3.961cm}}{\pgfqpoint{7.559cm}{-4.027cm}}{\pgfqpoint{7.487cm}{-4.087cm}}
\pgfpathcurveto{\pgfqpoint{7.417cm}{-4.145cm}}{\pgfqpoint{7.332cm}{-4.174cm}}{\pgfqpoint{7.231cm}{-4.174cm}}
\pgfpathcurveto{\pgfqpoint{7.173cm}{-4.174cm}}{\pgfqpoint{7.122cm}{-4.165cm}}{\pgfqpoint{7.081cm}{-4.147cm}}
\pgfpathcurveto{\pgfqpoint{7.039cm}{-4.131cm}}{\pgfqpoint{7.004cm}{-4.107cm}}{\pgfqpoint{6.975cm}{-4.078cm}}
\pgfpathcurveto{\pgfqpoint{6.947cm}{-4.047cm}}{\pgfqpoint{6.916cm}{-4.001cm}}{\pgfqpoint{6.882cm}{-3.94cm}}
\pgfpathcurveto{\pgfqpoint{6.851cm}{-3.877cm}}{\pgfqpoint{6.823cm}{-3.822cm}}{\pgfqpoint{6.8cm}{-3.775cm}}
\pgfpathmoveto{\pgfqpoint{5.324cm}{-2.923cm}}
\pgfpathlineto{\pgfqpoint{6.423cm}{-2.923cm}}
\pgfpathlineto{\pgfqpoint{5.869cm}{-1.656cm}}
\pgfpathlineto{\pgfqpoint{5.324cm}{-2.923cm}}
\pgfusepath{fill}
\end{scope}

%%%M
\begin{scope}[xshift=1cm]
\pgfpathmoveto{\pgfqpoint{9.392cm}{-3.659cm}}
\pgfpathlineto{\pgfqpoint{8.802cm}{-1.704cm}}
\pgfpathlineto{\pgfqpoint{8.802cm}{-3.823cm}}
\pgfpathcurveto{\pgfqpoint{8.802cm}{-3.94cm}}{\pgfqpoint{8.77cm}{-4.028cm}}{\pgfqpoint{8.707cm}{-4.087cm}}
\pgfpathcurveto{\pgfqpoint{8.645cm}{-4.145cm}}{\pgfqpoint{8.562cm}{-4.174cm}}{\pgfqpoint{8.458cm}{-4.174cm}}
\pgfpathcurveto{\pgfqpoint{8.358cm}{-4.174cm}}{\pgfqpoint{8.276cm}{-4.146cm}}{\pgfqpoint{8.212cm}{-4.089cm}}
\pgfpathcurveto{\pgfqpoint{8.149cm}{-4.03cm}}{\pgfqpoint{8.117cm}{-3.941cm}}{\pgfqpoint{8.117cm}{-3.823cm}}
\pgfpathlineto{\pgfqpoint{8.117cm}{-1.394cm}}
\pgfpathcurveto{\pgfqpoint{8.117cm}{-1.26cm}}{\pgfqpoint{8.159cm}{-1.17cm}}{\pgfqpoint{8.243cm}{-1.124cm}}
\pgfpathcurveto{\pgfqpoint{8.326cm}{-1.077cm}}{\pgfqpoint{8.439cm}{-1.053cm}}{\pgfqpoint{8.581cm}{-1.053cm}}
\pgfpathlineto{\pgfqpoint{8.812cm}{-1.053cm}}
\pgfpathcurveto{\pgfqpoint{8.951cm}{-1.053cm}}{\pgfqpoint{9.052cm}{-1.063cm}}{\pgfqpoint{9.113cm}{-1.084cm}}
\pgfpathcurveto{\pgfqpoint{9.177cm}{-1.105cm}}{\pgfqpoint{9.224cm}{-1.143cm}}{\pgfqpoint{9.254cm}{-1.197cm}}
\pgfpathcurveto{\pgfqpoint{9.284cm}{-1.252cm}}{\pgfqpoint{9.318cm}{-1.34cm}}{\pgfqpoint{9.357cm}{-1.463cm}}
\pgfpathlineto{\pgfqpoint{9.891cm}{-3.144cm}}
\pgfpathlineto{\pgfqpoint{10.426cm}{-1.463cm}}
\pgfpathcurveto{\pgfqpoint{10.464cm}{-1.34cm}}{\pgfqpoint{10.499cm}{-1.252cm}}{\pgfqpoint{10.529cm}{-1.197cm}}
\pgfpathcurveto{\pgfqpoint{10.559cm}{-1.143cm}}{\pgfqpoint{10.605cm}{-1.105cm}}{\pgfqpoint{10.667cm}{-1.084cm}}
\pgfpathcurveto{\pgfqpoint{10.73cm}{-1.063cm}}{\pgfqpoint{10.832cm}{-1.053cm}}{\pgfqpoint{10.971cm}{-1.053cm}}
\pgfpathlineto{\pgfqpoint{11.201cm}{-1.053cm}}
\pgfpathcurveto{\pgfqpoint{11.344cm}{-1.053cm}}{\pgfqpoint{11.457cm}{-1.077cm}}{\pgfqpoint{11.54cm}{-1.124cm}}
\pgfpathcurveto{\pgfqpoint{11.624cm}{-1.17cm}}{\pgfqpoint{11.666cm}{-1.26cm}}{\pgfqpoint{11.666cm}{-1.394cm}}
\pgfpathlineto{\pgfqpoint{11.666cm}{-3.823cm}}
\pgfpathcurveto{\pgfqpoint{11.666cm}{-3.94cm}}{\pgfqpoint{11.634cm}{-4.028cm}}{\pgfqpoint{11.57cm}{-4.087cm}}
\pgfpathcurveto{\pgfqpoint{11.508cm}{-4.145cm}}{\pgfqpoint{11.425cm}{-4.174cm}}{\pgfqpoint{11.319cm}{-4.174cm}}
\pgfpathcurveto{\pgfqpoint{11.221cm}{-4.174cm}}{\pgfqpoint{11.14cm}{-4.145cm}}{\pgfqpoint{11.076cm}{-4.087cm}}
\pgfpathcurveto{\pgfqpoint{11.012cm}{-4.028cm}}{\pgfqpoint{10.981cm}{-3.94cm}}{\pgfqpoint{10.981cm}{-3.823cm}}
\pgfpathlineto{\pgfqpoint{10.981cm}{-1.704cm}}
\pgfpathlineto{\pgfqpoint{10.391cm}{-3.659cm}}
\pgfpathcurveto{\pgfqpoint{10.352cm}{-3.786cm}}{\pgfqpoint{10.321cm}{-3.88cm}}{\pgfqpoint{10.295cm}{-3.94cm}}
\pgfpathcurveto{\pgfqpoint{10.272cm}{-3.999cm}}{\pgfqpoint{10.228cm}{-4.052cm}}{\pgfqpoint{10.162cm}{-4.101cm}}
\pgfpathcurveto{\pgfqpoint{10.097cm}{-4.15cm}}{\pgfqpoint{10.007cm}{-4.174cm}}{\pgfqpoint{9.891cm}{-4.174cm}}
\pgfpathcurveto{\pgfqpoint{9.804cm}{-4.174cm}}{\pgfqpoint{9.731cm}{-4.158cm}}{\pgfqpoint{9.671cm}{-4.126cm}}
\pgfpathcurveto{\pgfqpoint{9.61cm}{-4.096cm}}{\pgfqpoint{9.563cm}{-4.056cm}}{\pgfqpoint{9.53cm}{-4.007cm}}
\pgfpathcurveto{\pgfqpoint{9.497cm}{-3.958cm}}{\pgfqpoint{9.47cm}{-3.904cm}}{\pgfqpoint{9.45cm}{-3.846cm}}
\pgfpathcurveto{\pgfqpoint{9.431cm}{-3.786cm}}{\pgfqpoint{9.412cm}{-3.724cm}}{\pgfqpoint{9.392cm}{-3.659cm}}
\pgfusepath{fill}
\end{scope}

%%A
\begin{scope}[xshift=3cm]
\pgfpathmoveto{\pgfqpoint{14.79cm}{-3.775cm}}
\pgfpathlineto{\pgfqpoint{14.614cm}{-3.389cm}}
\pgfpathlineto{\pgfqpoint{13.119cm}{-3.389cm}}
\pgfpathlineto{\pgfqpoint{12.943cm}{-3.783cm}}
\pgfpathcurveto{\pgfqpoint{12.874cm}{-3.937cm}}{\pgfqpoint{12.816cm}{-4.04cm}}{\pgfqpoint{12.767cm}{-4.095cm}}
\pgfpathcurveto{\pgfqpoint{12.719cm}{-4.148cm}}{\pgfqpoint{12.639cm}{-4.174cm}}{\pgfqpoint{12.529cm}{-4.174cm}}
\pgfpathcurveto{\pgfqpoint{12.435cm}{-4.174cm}}{\pgfqpoint{12.352cm}{-4.146cm}}{\pgfqpoint{12.281cm}{-4.089cm}}
\pgfpathcurveto{\pgfqpoint{12.209cm}{-4.031cm}}{\pgfqpoint{12.173cm}{-3.967cm}}{\pgfqpoint{12.173cm}{-3.894cm}}
\pgfpathcurveto{\pgfqpoint{12.173cm}{-3.852cm}}{\pgfqpoint{12.181cm}{-3.809cm}}{\pgfqpoint{12.198cm}{-3.764cm}}
\pgfpathcurveto{\pgfqpoint{12.214cm}{-3.719cm}}{\pgfqpoint{12.242cm}{-3.657cm}}{\pgfqpoint{12.281cm}{-3.578cm}}
\pgfpathlineto{\pgfqpoint{13.222cm}{-1.585cm}}
\pgfpathcurveto{\pgfqpoint{13.248cm}{-1.527cm}}{\pgfqpoint{13.28cm}{-1.459cm}}{\pgfqpoint{13.317cm}{-1.38cm}}
\pgfpathcurveto{\pgfqpoint{13.355cm}{-1.299cm}}{\pgfqpoint{13.396cm}{-1.232cm}}{\pgfqpoint{13.437cm}{-1.179cm}}
\pgfpathcurveto{\pgfqpoint{13.481cm}{-1.126cm}}{\pgfqpoint{13.537cm}{-1.083cm}}{\pgfqpoint{13.606cm}{-1.051cm}}
\pgfpathcurveto{\pgfqpoint{13.676cm}{-1.017cm}}{\pgfqpoint{13.762cm}{-1.001cm}}{\pgfqpoint{13.864cm}{-1.001cm}}
\pgfpathcurveto{\pgfqpoint{13.968cm}{-1.001cm}}{\pgfqpoint{14.054cm}{-1.017cm}}{\pgfqpoint{14.123cm}{-1.051cm}}
\pgfpathcurveto{\pgfqpoint{14.193cm}{-1.083cm}}{\pgfqpoint{14.249cm}{-1.125cm}}{\pgfqpoint{14.291cm}{-1.176cm}}
\pgfpathcurveto{\pgfqpoint{14.334cm}{-1.228cm}}{\pgfqpoint{14.37cm}{-1.284cm}}{\pgfqpoint{14.399cm}{-1.344cm}}
\pgfpathcurveto{\pgfqpoint{14.429cm}{-1.403cm}}{\pgfqpoint{14.466cm}{-1.481cm}}{\pgfqpoint{14.512cm}{-1.581cm}}
\pgfpathlineto{\pgfqpoint{15.473cm}{-3.561cm}}
\pgfpathcurveto{\pgfqpoint{15.548cm}{-3.712cm}}{\pgfqpoint{15.586cm}{-3.821cm}}{\pgfqpoint{15.586cm}{-3.89cm}}
\pgfpathcurveto{\pgfqpoint{15.586cm}{-3.961cm}}{\pgfqpoint{15.55cm}{-4.027cm}}{\pgfqpoint{15.478cm}{-4.087cm}}
\pgfpathcurveto{\pgfqpoint{15.408cm}{-4.145cm}}{\pgfqpoint{15.322cm}{-4.174cm}}{\pgfqpoint{15.222cm}{-4.174cm}}
\pgfpathcurveto{\pgfqpoint{15.163cm}{-4.174cm}}{\pgfqpoint{15.113cm}{-4.165cm}}{\pgfqpoint{15.071cm}{-4.147cm}}
\pgfpathcurveto{\pgfqpoint{15.029cm}{-4.131cm}}{\pgfqpoint{14.994cm}{-4.107cm}}{\pgfqpoint{14.966cm}{-4.078cm}}
\pgfpathcurveto{\pgfqpoint{14.937cm}{-4.047cm}}{\pgfqpoint{14.906cm}{-4.001cm}}{\pgfqpoint{14.873cm}{-3.94cm}}
\pgfpathcurveto{\pgfqpoint{14.841cm}{-3.877cm}}{\pgfqpoint{14.814cm}{-3.822cm}}{\pgfqpoint{14.79cm}{-3.775cm}}
\pgfpathmoveto{\pgfqpoint{13.314cm}{-2.923cm}}
\pgfpathlineto{\pgfqpoint{14.414cm}{-2.923cm}}
\pgfpathlineto{\pgfqpoint{13.859cm}{-1.656cm}}
\pgfpathlineto{\pgfqpoint{13.314cm}{-2.923cm}}
\pgfusepath{fill}
\end{scope}

\draw[draw=none,fill=textLogoColor] (-5,-0.5)--(22,-0.5) %%..controls (7,-0.5) and (10,-0.5)..(17,-0.5)
.. controls(17,-0.8) and (10,-0.9)..(8.5,-0.9)..controls(6,-0.9) and (0,-0.8)..(-5,-0.5);

\draw[text=textLogoColor!80!white](8.5,-3.5)node[anchor=north]{\sffamily\fontsize{6}{6}\selectfont Laboratoire de Math\'ematiques};
\draw[text=textLogoColor!80!white](8.5,-5.5)node[anchor=north]{\sffamily\fontsize{6}{6}\selectfont Universit\'e de Savoie};

  
\end{scope}



   %%       \end{tikzpicture}
   \small
   \begin{tabular}{cc}
     \begin{minipage}{0.48\textwidth}
       \includegraphics[width=1.8cm]{liris-logo}
       \quad
       \begin{minipage}{0.55\textwidth}
         D. C{\oe}urjolly\\
         G. Damiand\\
         T. Roussillon\\
       \end{minipage}
     \end{minipage} &
     \begin{minipage}{0.48\textwidth}
       \includegraphics[width=1.8cm]{lama-logo}
       \quad
       \begin{minipage}{0.55\textwidth}
         J.-O. Lachaud\\[1cm]
       \end{minipage}
     \end{minipage} \\
     \begin{minipage}{0.48\textwidth}
       \includegraphics[width=1.8cm]{loria-logo_new}
       \quad
       \begin{minipage}{0.55\textwidth}
         B. Kerautret\\[1cm]
       \end{minipage}
     \end{minipage} &
     \begin{minipage}{0.48\textwidth}
       \includegraphics[width=1.8cm]{greyc-logo}
       \quad
       \begin{minipage}{0.55\textwidth}
         S. Fourey\\[0.1cm]
       \end{minipage}
     \end{minipage} \\
     \begin{minipage}{0.48\textwidth}
       \includegraphics[width=1.8cm]{gipsa-logo}
       \quad
       \begin{minipage}{0.55\textwidth}
         I. Sivignon\\[1cm]
       \end{minipage}
     \end{minipage}&\\
   \end{tabular}

\end{frame}
%------------------------------------------------------------------------------

%------------------------------------------------------------------------------



%% \begin{itemize}
%%   \item Arbitrary adjacencies in $\Z^n$, but also in subdomains
%%   \item Digital topology = couple of adjacencies (Rosenfeld)
%%   \item Object = Topology + Set
%%   \item Operations: neighborhoods, border, connectedness and connected
%%     components, decomposition into digital layers, simple points
%%   \end{itemize}

%% \end{frame}
%% %------------------------------------------------------------------------------

%% %------------------------------------------------------------------------------
%% \begin{frame}
%%   \frametitle{Adjacency}

%%   \alert{Genericity} $\Rightarrow$ concept \alert{\texttt{CAdjacency}}

%%   \begin{itemize}
%%   \item Types: \texttt{Space}, \texttt{Point}, \texttt{Adjacency}
%%   \item Methods: 
%%     \begin{itemize}
%%     \item \texttt{isAdjacentTo}( p1, p2 )
%%     \item \texttt{isProperlyAdjacentTo}( p1, p2 )
%%     \item \texttt{writeNeighborhood}( p, outit )
%%     \item \texttt{writeProperNeighborhood}( p, outit )
%%     \item \texttt{writeNeighborhood}( p, outit, pred )
%%     \item \texttt{writeProperNeighborhood}( p, outit, pred )
%%     \end{itemize}
%%   \item Models: 
%%     \begin{itemize}
%%     \item \texttt{MetricAdjacency}: 4-, 8-, 6-, 18-, 26-, $2n$-, $3^n-1$- adjacencies
%%     \item \texttt{DomainAdjacency}: adjacency limited by a specified domain.
%%     \end{itemize}
%%   \end{itemize}
%% \end{frame}
%% %------------------------------------------------------------------------------

%% %------------------------------------------------------------------------------
%% \begin{frame}[fragile]
%%   \frametitle{Digital topology}
%%   \alert{Digital topology} = couple of instances of adjacencies

%%   \begin{itemize}
%%   \item template class \texttt{DigitalTopology}
%%     \scriptsize
%%     \begin{semiverbatim}
%%  typedef SpaceND< 3,int > Z3;
%%  typedef MetricAdjacency< Z3, 1 > Adj6;
%%  typedef MetricAdjacency< Z3, 2 > Adj18;
%%  typedef DigitalTopology< Adj6, Adj18 > DT6_18;

%%  Adj6 adj6;
%%  Adj18 adj18;
%%  DT6_18 dt6_18( adj6, adj18, JORDAN_DT );
%%     \end{semiverbatim}
%%     \normalsize
%%   \item Jordan topologies may be specified (for future use)
%%   \item instances are necessary (e.g., adj may not be invariant by translation)
%%   \item reverse topology is the reversed couple
%%   \end{itemize}
%% \end{frame}
%% %------------------------------------------------------------------------------

%% %------------------------------------------------------------------------------
%% \begin{frame}[fragile]
%%   \frametitle{Digital Object}

%%   \alert{Digital object} = topology $+$ digital set

%%   \begin{itemize}
%%   \item template class \texttt{Object}

%%     \scriptsize
%%     \begin{semiverbatim}
%%  typedef HyperRectDomain< Z3 > Domain; 
%%  typedef DigitalSetSelector<Domain, BIG_DS+HIGH_BEL_DS>::Type DigitalSet;
%%  typedef Object<DT6_18, DigitalSet> ObjectType;
%%  Point p1( -50, -50, -50 ); Point p2( 50, 50, 50 );
%%  Domain domain( p1, p2 );
%%  // ball of radius 30
%%  DigitalSet ball_set( domain );
%%  Shapes<Domain>::addNorm2Ball( ball_set, Point( 0, 0 ), 30 );
%%  ObjectType ball_object( dt6_18, ball_set );
%%     \end{semiverbatim}
%%     \normalsize
%%   \item Objects may be passed by value and copied without cost
%%   \item Methods:
%%     \begin{itemize}
%%     \item neighborhoods, border, geodesic neighborhoods are objects
%%     \item (lazy) connectedness, connected components
%%     \item simple points (in Z2 and Z3)
%%     \end{itemize}
%%   \end{itemize}

%% \end{frame}
%% %------------------------------------------------------------------------------

%% %------------------------------------------------------------------------------
%% \begin{frame}[fragile]
%%   \frametitle{Expander: digital layers in an object}
  
%%   \begin{itemize}
%%   \item Expansion layer by layer within an object, starting from an initial core
%%   \item core = a point or a pointset specified by iterators
%%   \item each new layer = the set of points of the object adjacent to
%%     the preceding layer
%%   \item each layer is iterable, has a digital distance to core
%%   \item finished when no more neighbor expansion is possible
%%   \item useful for \alert{connectedness}
%%   \end{itemize}
%% \end{frame}
%% %------------------------------------------------------------------------------

%% %------------------------------------------------------------------------------
%% \begin{frame}
%%   \frametitle{Expander: digital layers in an object}
%%   \begin{center}
%%     \begin{tabular}{cc}
%%       \includegraphics[width=0.4\textwidth]{house-layers4-4} &
%%       \includegraphics[width=0.4\textwidth]{house-layers4-8} \\
%%       background in 4-adj &
%%       background in 8-adj \\
%%     \end{tabular}

%%     \texttt{tests/topology/testSimpleExpander.cpp}
%%   \end{center}
%% \end{frame}

%% %------------------------------------------------------------------------------
%% \begin{frame}[fragile]
%%   \frametitle{Example: greedy homotopic thinning}
%%     \scriptsize
%%     \begin{semiverbatim}
%%   int layer = 0;
%%   do \{
%%       DigitalSet & S = shape.pointSet();
%%       std::queue<DigitalSet::Iterator> Q;
%%       for ( DigitalSet::Iterator it = S.begin(); it != S.end(); ++it )
%%         if ( shape.\alertred{isSimple}( *it ) )
%%           Q.push( it );
%%       nb_simple = 0;
%%       while ( ! Q.empty() ) \{
%%         DigitalSet::Iterator it = Q.front();
%%         Q.pop();
%%         if ( shape.isSimple( *it ) ) \{
%%           S.erase( *it );
%%           ++nb_simple;
%%         \}
%%       \}
%%       ++layer;
%%   \} while ( nb_simple != 0 );
%%     \end{semiverbatim}
%%     \normalsize
%%     See \texttt{testObject.cpp}
%% \end{frame}
%% %------------------------------------------------------------------------------

%% %------------------------------------------------------------------------------
%% \begin{frame}[fragile]
%%   \frametitle{Example: greedy homotopic thinning}
%%   \begin{center}
%%     \begin{tabular}{cc}
%%       \includegraphics[width=0.4\textwidth]{shape-thinning-4-8} &
%%       \includegraphics[width=0.4\textwidth]{shape-thinning-8-4} \\
%%       thinning in (4,8) &
%%       thinning in (8,4) \\
%%     \end{tabular}

%%     \texttt{tests/topology/testObject.cpp}
%%   \end{center}
%% \end{frame}
%% %------------------------------------------------------------------------------

%% %------------------------------------------------------------------------------
%% \begin{frame}
%%   \frametitle{Conclusion and perspectives}
  
%%   \begin{itemize}
%%   \item complete Rosenfeld's approach: curves and separation
%%   \item whole digital topology framework of Herman and Udupa
%%     \begin{itemize}
%%     \item digital surface as a couple of $\omega$-adjacent points
%%     \item immediate interior and exterior, interior and exterior
%%     \item $\kappa \lambda$-borders, $\kappa \lambda$-boundaries
%%     \item digital pictures
%%     \end{itemize}
%%   \item interpixel topology or cartesian cellular grid topology 
%%   \end{itemize}
%%   See on-line doc.: \Cb{Digital topology and digital objects}

%% \end{frame}
%% %------------------------------------------------------------------------------

\end{document}
